% =====================================================================
% CHAPTER 1: INTRODUCTION                                   budget: 3 pages
% ROLE: Say why the question matters, state Q1-Q3, and preview the answer.
% Write this chapter last, when the results chapters are fixed.
% =====================================================================
\chapter{Introduction}
\label{ch:intro}

\section{Motivation}

% P1 KEY: Communication inside an organization shapes its decisions, and
%         hierarchy shapes that communication.
% SUPPORT:
%   - Exposé p.2: communication drives decision efficiency and alignment.
%   - Email is the main written channel in firms. Enron is the standard
%     public example.
%   - Power shows in workplace language (Gilbert 2012; Prabhakaran et al.).

% P2 KEY: LLMs are now used as stand-ins for people (as simulated agents)
%         and as evaluators of text (as judges).
% SUPPORT:
%   - Agents: Horton et al. ("Homo silicus"); Park et al. 2023 (Generative Agents).
%   - Judges: Zheng et al. 2023 (MT-Bench, LLM-as-a-judge).
%   - Promise: controlled experiments on role and hierarchy that are hard
%     to run with real employees.

% P3 KEY: Whether this works for organizational email is open: do simulated
%         replies show real hierarchy patterns, and can a judge score them?
% SUPPORT:
%   - Power-asymmetric conversations paper: LLM agents mirror some
%     socio-cognitive effects, not all.
%   - Moral Mazes (HR Simulator): LLM email tone differs from human email.
%   - Few studies compare simulated email to real email on the same measures.

\section{Research questions}

% P4 KEY: This thesis builds an email simulator with role personas and
%         asks three questions.
% SUPPORT:
%   - Q1: How does hierarchical role influence response content, sentiment
%     and decision attitude?
%   - Q2: Are LLM-generated responses diverse enough and similar to
%     empirical patterns?
%   - Q3: Can LLM-as-a-judge give consistent and calibrated scores for
%     communication quality?
%   - Wording from the Exposé (p.3). Keep it.

% P5 KEY: How each question is made measurable.
% SUPPORT:
%   - Q1: writing direction (up / to a peer / down) as the treatment.
%     Outcomes: orders, tone, politeness, decision field.
%   - Q2: generated vs. real replies to the same Enron emails, plus Self-BLEU.
%   - Q3: self-agreement, agreement between judges, self-preference, effect
%     of showing the incoming message, check against people.

% P6 KEY: The thesis uses only free models, which makes it a study of what
%         openly available LLMs can do.
% SUPPORT:
%   - No paid API (agreed with the supervisor).
%   - 4 models: Llama 3.2 3B (local), gpt-oss-20b, gpt-oss-120b, qwen3.8-27b
%     (free hosted tiers). Total cost $0.
%   - SRC: RO 0 (Models).

\section{Contributions}

% P7 KEY: What this thesis adds.
% SUPPORT:
%   - A reproducible simulator: 4 models, 3 designs, about 12,000 cached replies.
%   - A "real email first" rule: every hierarchy measure is tested on
%     real Enron email before it is used on generated replies.
%   - Q1: the hierarchy effect exists but is fragile (personas, scenario, model).
%   - Q2: the gap to real email is style, not topic.
%   - Q3: a corrected self-preference test (neutral reference writer) and
%     evidence that the judge needs the incoming message.

% P8 KEY: The answer in short (the main story from main.tex, in 3-4 sentences).
% SUPPORT:
%   - Copy the MAIN STORY block and shorten it once chapters 5-7 are written.

\section{Outline}

% P9 KEY: One sentence per chapter.
% SUPPORT:
%   - Ch. 2 background, Ch. 3 data and simulator, Ch. 4 measures,
%     Ch. 5-7 one chapter per question, Ch. 8 discussion and conclusion.
